% Generated draft: requires mathematical and source review.
\documentclass[11pt]{article}
\usepackage[margin=25mm]{geometry}
\usepackage{fontspec}
\setmainfont{Latin Modern Roman}
\usepackage{amsmath,amssymb,mathtools,unicode-math}
\setmathfont{Latin Modern Math}
\usepackage{enumitem,needspace,xurl,hyperref}
\setlength{\parindent}{0pt}
\setlength{\parskip}{5pt}
\setlength{\emergencystretch}{3em}
\newcommand{\N}{\mathbb{N}}
\newcommand{\Z}{\mathbb{Z}}
\newcommand{\Q}{\mathbb{Q}}
\newcommand{\R}{\mathbb{R}}
\newcommand{\C}{\mathbb{C}}
\newcommand{\F}{\mathbb{F}}
\newcommand{\A}{\mathbb{A}}
\newcommand{\PP}{\mathbb{P}}
\newcommand{\E}{\mathbb{E}}
\newcommand{\eps}{\varepsilon}
\newcommand{\dd}{\,\mathrm{d}}
\begin{document}
\section*{UnsolvedMath ID: 10800002}
\textbf{Source problem code:} catalogue entry
\subsection{Definitions and mathematical statement}
Let \(f\) be the given nonsimple singularity and \(f_\lambda\) a morsification in the deformation family intended by the source. In the usual complex isolated-hypersurface setting, a morsification has nondegenerate critical points and suitable critical values, and a distinguished system of paths to those values specifies vanishing cycles and their Dynkin diagram \(D\). The record does not state that full setting, the deformation family, or the diagram conventions. Moreover, its wording says paths from \(0\) to ``critical points,'' whereas the usual distinguished paths lie in the value plane and end at critical values; this distinction remains unresolved.

Subject to recovering that setup, fix the critical-value data \(\mathbf v\) of \(f_\lambda\) and the specified path system \(\gamma\). Let
\[
 \mathcal F(D)=\{\beta\cdot D:\beta\text{ is a braid acting by formal Picard--Lefschetz moves}\},
\]
and let
\[
 \mathcal R(\mathbf v,\gamma)
 =\{D(f_{\lambda'},\gamma): f_{\lambda'}\text{ is an allowed morsification of }f
 \text{ with the same critical values }\mathbf v\}.
\]
The allowed parameter domain, braid action, ordering conventions, and equality of graphs must be fixed by the omitted source context.

\emph{Problem (Victor Vassiliev, Problem 1B).} Determine which Dynkin graphs are formally obtainable from \(D\) by a braid-defined sequence of Picard--Lefschetz moves but cannot occur for a morsification with the same critical values and the same path system; that is, characterize
\[
 \mathcal F(D)\setminus\mathcal R(\mathbf v,\gamma).
\]
The path system and critical-value constraint are not discarded after a formal move. The question asks for a characterization of the difference, not an unsupported assertion that it is nonempty for every nonsimple singularity.
\subsection{Short English statement}
For a fixed nonsimple singularity, identify the Dynkin graphs reached by formal braid/Picard--Lefschetz operations that cannot be realized by actual morsifications with the prescribed critical values and path system. The deformation and path conventions need recovery from the omitted setup.
\subsection{Sources}
\begin{itemize}
\item Local UnsolvedMath registry, permanent ID 10800002, source code catalogue entry; proposer Victor Vassiliev, Problem 1B.
\item Vassiliev, \emph{A Few Problems on Monodromy and Discriminants}, Problem 1B. Supplied URL: \url{https://amj.math.stonybrook.edu/html-articles/Files-2015-2024/15-11/}; recorded access 2026-07-29; extraction html-structure. The article contains eighteen explicitly labeled theorem-style problem environments.
\item Supplied status open and OPEN-TRIAGE. The research summary says formal and geometric reachable sets agree for simple singularities and diverge in general for nonsimple ones, while reporting no published explicit counterexample graphs found. These statements and the problem's apparent continued openness are editorial assessments without separate citations in this record. Current primary-literature status and redistribution terms for the publicly accessible author/publisher copy remain NEEDS\_REVIEW.
\end{itemize}
\end{document}
